\begin{table}[htbp]
\centering
\caption{逻辑回归+边界SMOTE高质量类TP/FP/FN特征均值对比}
\label{tab:high-quality-error-features}
\small
\setlength{\tabcolsep}{4pt}
\resizebox{\textwidth}{!}{%
\begin{tabular}{lccccc}
\toprule
特征 & TP均值 & FP均值 & FN均值 & FP-TP标准化差异 & FN-TP标准化差异 \\
\midrule
出勤率 & 88.794 & 87.771 & 99.890 & -0.088 & 0.959 \\
作业完成率 & 77.504 & 71.485 & 52.560 & -0.419 & -1.737 \\
平均作业分数 & 90.269 & 88.537 & 79.820 & -0.120 & -0.724 \\
考试分数 & 90.880 & 86.437 & 98.710 & -0.262 & 0.461 \\
平均测验分数 & 68.694 & 69.952 & 62.380 & 0.073 & -0.369 \\
参与频率 & 13.921 & 13.235 & 9.000 & -0.081 & -0.578 \\
学习管理系统登录次数 & 29.842 & 30.059 & 9.000 & 0.017 & -1.607 \\
平均会话时长 & 36.750 & 36.498 & 47.310 & -0.018 & 0.737 \\
视频完成率 & 59.709 & 56.772 & 75.750 & -0.147 & 0.801 \\
自我评估分数 & 4.018 & 3.809 & 4.610 & -0.185 & 0.522 \\
同行反馈分数 & 2.993 & 2.995 & 3.960 & 0.002 & 0.831 \\
教师反馈分数 & 4.111 & 3.872 & 2.240 & -0.206 & -1.614 \\
参与度指数 & 0.544 & 0.527 & 0.400 & -0.118 & -0.992 \\
\bottomrule
\end{tabular}
}%
\vspace{2pt}
\parbox{\textwidth}{\footnotesize 注：TP表示真实高质量且预测为高质量，FP表示真实非高质量但预测为高质量，FN表示真实高质量但预测为非高质量。标准化差异按（组均值-TP均值）/全样本标准差计算。}
\end{table}
